\documentclass[11pt]{article}
\usepackage[T1]{fontenc}
\usepackage[utf8]{inputenc}
\usepackage{lmodern}
\usepackage{amsmath,amssymb,amsthm}
\usepackage[margin=1in]{geometry}
\usepackage{microtype}
\usepackage[colorlinks=true,linkcolor=blue,citecolor=blue,urlcolor=blue]{hyperref}
\hypersetup{
  pdftitle={An infinite family of counterexamples to Stanton's polynomial positivity conjecture},
  pdfauthor={Zizhou Zhang},
  pdfsubject={Stanton's 1999 Conjecture 12},
  pdfkeywords={coefficient positivity, cyclotomic polynomials, integer partitions, counterexamples, AI construction}
}

\newtheorem{theorem}{Theorem}
\newcommand{\NN}{\mathbb Z_{\geq0}}

\title{An infinite family of counterexamples to\newline
Stanton's polynomial positivity conjecture}
\author{Zizhou Zhang\\[0.5em]
\small Academy of Mathematics and Systems Science\\
\small Chinese Academy of Sciences\\
\small Beijing, China}
\date{November 26, 2026}

\begin{document}
\maketitle

\begin{abstract}
Conjecture~12 in Stanton's 1999 survey asserts that a certain factorial
product quotient associated with a positive integer sequence has
nonnegative coefficients whenever it is a polynomial. We present an
explicit counterexample of degree~816 and an infinite family extending
it. The member indexed by $r\geq0$ has degree $890r+816$ and coefficient
$-1$ at $q^{300r+277}$. Polynomiality follows from cyclotomic
factorization, and the negative coefficient is established by parity
and congruence arguments. The explicit example and the family were
constructed entirely by OpenAI's Codex (GPT-6).
\end{abstract}

\medskip
\noindent\textbf{Keywords.} Coefficient positivity; cyclotomic polynomials;
integer partitions; counterexamples.

\section{The conjecture and the example}

Let $a_1,\ldots,a_n$ be positive integers and put $N=\sum_{i=1}^n a_i$.
Stanton \cite[Conjecture~12]{Stanton} conjectured that, if
\begin{equation}\label{eq:stanton}
 G_{\boldsymbol a}(q)
 =\frac{\prod_{j=1}^{N}(1-q^j)}
        {\prod_{i=1}^{n}(1-q^i)^{a_i}}
\end{equation}
is a polynomial in $q$, then all of its coefficients are nonnegative.
Here the sequence length is denoted consistently by $n$.

Consider the rational function
\begin{equation}\label{eq:example}
 Q(q)=\frac{(1-q^{277})(1-q^{278})(1-q^{279})(1-q^{280})}
 {(1-q^{10})(1-q^{56})(1-q^{93})(1-q^{139})}.
\end{equation}

\begin{theorem}\label{thm:counterexample}
The function $Q(q)$ is an integer polynomial of degree $816$, and
\[
 [q^{277}]Q(q)=-1.
\]
It is an instance of \eqref{eq:stanton} with $n=276$ and $N=280$.
Consequently, Stanton's Conjecture~12 is false.
\end{theorem}

\begin{proof}
We first prove polynomiality. Set
\[
 A=\{277,278,279,280\},\qquad B=\{10,56,93,139\},
\]
and, for a finite set $E$ and an integer $d\geq1$, write
$\nu_d(E)=\#\{e\in E:d\mid e\}$. By the cyclotomic factorization
$q^t-1=\prod_{d\mid t}\Phi_d(q)$, the exponent of $\Phi_d(q)$ in
\eqref{eq:example} is
\[
 e_d=\nu_d(A)-\nu_d(B).
\]
For $d=1$, both counts are four. For $d=2$, both counts are two:
the even elements of $A$ are $278,280$, and those of $B$ are $10,56$.
The only pair of distinct elements of $B$ with a common divisor greater
than one is $10,56$, whose greatest common divisor is~$2$.
Thus, for $d>2$, at most one element of $B$ is divisible by $d$.
Whenever such an element exists, a multiple of it occurs in $A$, since
\[
 10\mid280,\qquad 56\mid280,\qquad 93\mid279,\qquad 139\mid278.
\]
It follows that $e_d\geq0$ for every $d\geq1$, proving
$Q(q)\in\mathbb Z[q]$. Its degree is
\[
 277+278+279+280-(10+56+93+139)=816.
\]

We next compute the coefficient at degree $277$. Define $p(t)$ by the
formal power series
\[
 H(q)=\frac1{(1-q^{10})(1-q^{56})(1-q^{93})(1-q^{139})}
     =\sum_{t\geq0}p(t)q^t.
\]
Then $p(t)$ counts the solutions in $\NN^4$ of
\[
 10x+56y+93z+139w=t.
\]
We claim that $p(277)=0$. Indeed, a solution at $t=277$ would have
$w\in\{0,1\}$, because $2\cdot139=278$, and $z+w$ would be odd.
If $w=1$, the remaining total is $138$. Hence $z=0$, and
\[
 5x+28y=69.
\]
This gives $0\leq y\leq2$, whereas reduction modulo $5$ gives
$y\equiv3\pmod5$, a contradiction. If $w=0$, then $z$ is odd and
$93z\leq277<3\cdot93$, so $z=1$. We obtain
\[
 5x+28y=92.
\]
Now $0\leq y\leq3$, whereas $y\equiv4\pmod5$, another contradiction.
Therefore $p(277)=0$.

The numerator of \eqref{eq:example} has constant coefficient one,
coefficient $-1$ at degree $277$, and no terms of degrees $1$ through
$276$. Since $p(0)=1$, it follows that
\[
 [q^{277}]Q(q)=p(277)-p(0)=-1.
\]

Finally, take $n=276$ and set
\[
 a_i=\begin{cases}
 2,&i\in\{10,56,93,139\},\\
 1,&\text{otherwise},\quad 1\leq i\leq276.
 \end{cases}
\]
All entries are positive, and their sum is $N=280$.
Cancelling one copy of each factor $1-q^i$, for $1\leq i\leq276$,
in \eqref{eq:stanton} gives exactly $Q(q)$. Thus the example satisfies
every hypothesis of the conjecture and violates its conclusion.
\end{proof}

\section{An infinite family}

The same congruence obstruction gives a counterexample for every
nonnegative integer $r$. Define
\begin{equation}\label{eq:family}
 Q_r(q)=
 \frac{\prod_{j=277}^{280}(1-q^{300r+j})}
 {(1-q^{10})(1-q^{60r+56})(1-q^{100r+93})(1-q^{150r+139})}.
\end{equation}
The polynomial in \eqref{eq:example} is $Q_0(q)$.

\begin{theorem}\label{thm:family}
For every integer $r\geq0$, the function $Q_r(q)$ belongs to
$\mathbb Z[q]$ and satisfies
\[
 \deg Q_r=890r+816,
 \qquad [q^{300r+277}]Q_r(q)=-1.
\]
Each $Q_r$ is an instance of \eqref{eq:stanton} with
$n=300r+276$ and $N=300r+280$. Consequently, these are infinitely many
distinct counterexamples to Stanton's Conjecture~12.
\end{theorem}

\begin{proof}
Fix $r\geq0$, and put
\[
 m=30r+28,\qquad u=5m-1,\qquad v=\frac{10m-1}{3},
 \qquad R=10m-3.
\]
Then \eqref{eq:family} becomes
\[
 Q_r(q)=\frac{\prod_{j=R}^{R+3}(1-q^j)}
 {(1-q^{10})(1-q^{2m})(1-q^v)(1-q^u)}.
\]
The integer $m$ is even, is congruent to $1$ modulo $3$, and is
congruent to $3$ modulo $5$. In particular, $u$ and $v$ are odd
integers, and
\[
 10<2m<v<u<R.
\]

We verify polynomiality as in Theorem~\ref{thm:counterexample}.
The four denominator weights have
$\gcd(10,2m)=2$, and every other pair is coprime.
Indeed, $u=5m-1$ implies $\gcd(m,u)=\gcd(5,u)=1$.
The identity $3v=10m-1$ implies $\gcd(m,v)=1$, and
$v=100r+93\equiv3\pmod5$ gives $\gcd(5,v)=1$.
Oddness of $u,v$ accounts for the factor $2$, and
\[
 3v-2u=1
\]
gives $\gcd(u,v)=1$. For $d=1$, the numerator and denominator each
have four exponents divisible by $d$; for $d=2$, each has two.
For $d>2$, at most one denominator weight is divisible by $d$.
Whenever it is, there is a numerator exponent divisible by that weight,
because
\[
 2u=R+1,\qquad 3v=R+2,\qquad 10\mid R+3,
 \qquad 2m\mid R+3.
\]
Thus every cyclotomic exponent is nonnegative, and $Q_r(q)\in\mathbb Z[q]$.

To determine the coefficient at degree $R$, suppose that
\begin{equation}\label{eq:family-representation}
 R=10x+2my+vz+uw,\qquad x,y,z,w\in\NN.
\end{equation}
Since $2u=R+1$ and $3v=R+2$, we have $w\leq1$ and $z\leq2$.
Also $z+w$ is odd. If $w=1$, then $z$ is even, and
\[
 u+2v-R=\frac{5m+4}{3}>0
\]
excludes $z=2$. Hence $z=0$, and \eqref{eq:family-representation}
reduces to
\[
 5x+my=\frac{5m}{2}-1.
\]
This implies $y\leq2$. On the other hand, reduction modulo $5$ gives
$3y\equiv4\pmod5$, hence $y\equiv3\pmod5$, a contradiction.

If $w=0$, then $z=1$, and \eqref{eq:family-representation} reduces to
\[
 5x+my=\frac{10m-4}{3}.
\]
Now $y\leq3$. Multiplying by $3$ and reducing modulo $5$ gives
$4y\equiv1\pmod5$, hence $y\equiv4\pmod5$, again a contradiction.
There is therefore no representation \eqref{eq:family-representation}.
The reciprocal denominator has coefficient zero at degree $R$ and
constant coefficient one. Since the numerator has coefficient $-1$
at degree $R$ and no terms of degrees $1,\ldots,R-1$, we obtain
\[
 [q^R]Q_r(q)=-1.
\]

The degree is
\[
 \deg Q_r=(40m-6)-(10+2m+u+v)
         =\frac{89m-44}{3}=890r+816.
\]
Finally, take $n=10m-4$ and let $a_i=2$ for
$i\in\{10,2m,v,u\}$ and $a_i=1$ for every other $1\leq i\leq n$.
All four exceptional indices lie in this range and are distinct, so
$\sum_{i=1}^n a_i=n+4=10m=N$. Cancelling the factors with indices
$1,\ldots,n$ in \eqref{eq:stanton} gives $Q_r$.
The strictly increasing degrees show that these counterexamples are
distinct.
\end{proof}

Both proofs use cyclotomic factorization and elementary congruences;
no polynomial expansion or computational search is required to verify
the results. The assertion addressed here is
Conjecture~12 of \cite{Stanton}. It is distinct from Stanton's fake
Gaussian sequence conjecture (Conjecture~11 in the same paper) and
from the positivity conjecture for quotients of Gaussian binomial
coefficients with the same upper argument.

\section*{Construction provenance and AI assistance}
The explicit counterexample and the infinite family of counterexamples
presented in this paper were constructed entirely by OpenAI's Codex
(GPT-6). Codex also assisted with proof development, exploratory
computations, manuscript drafting, and automated internal consistency
checks. The author takes responsibility for the content of this paper.

\begin{thebibliography}{1}
\bibitem{Stanton}
D.~Stanton,
\emph{Open positivity conjectures for integer partitions},
\emph{Trends in Mathematics} \textbf{2} (1999), 19--25.
Author's manuscript:
\url{https://www-users.cse.umn.edu/~stant001/PAPERS/korea99.pdf}.
\end{thebibliography}

\end{document}
